% =============================================================================
% CITA FRAMEWORK - CITA Paper in ALKALI Format
% =============================================================================

\vspace{-2mm}
\section{CITA: Contrastive Instruction-Tuned Alignment}
\label{sec:cita_framework}

Building on the unified loss formulation (Section~\ref{sec:unified_loss}), we present the complete CITA framework, including the training triplet structure, the full CITA-KL loss, and the training pipeline.

\vspace{-2mm}
\subsection{Training Triplet}
\label{sec:training_triplet}

Each CITA training example consists of a 4-tuple:

\begin{equation}
(I, X, Y^+, Y^-) \in \mathcal{D}_{\text{CITA}}
\label{eq:triplet}
\end{equation}

where:
\begin{itemize}[leftmargin=*,itemsep=1pt]
    \item $I$: Alignment instruction (e.g., ``Respond with safety-first principles'')
    \item $X$: User input/query
    \item $Y^+$: Preferred response (aligned with instruction $I$)
    \item $Y^-$: Rejected response (misaligned with instruction $I$)
\end{itemize}

\textbf{Crucial Difference from DPO.} In standard DPO~\cite{rafailov2023direct}, the instruction $I$ is absent or implicit. CITA \textit{explicitly conditions} preference on $I$, enabling the model to learn instruction-specific behavioral patterns---analogous to how contrastive learning~\cite{chen2020simple,khosla2020supervised} conditions representations on class labels.

\vspace{-2mm}
\subsection{CITA-KL Loss}
\label{sec:cita_kl_loss}

The complete CITA objective combines contrastive learning~\cite{oord2018representation,gao2021simcse,radford2021learning} with mandatory KL~\cite{korbak2022rl}:

\begin{equation}
\mathcal{L}_{\text{CITA-KL}} = -\log P^+ + \lambda \cdot \text{KL}[\pi_\theta \,\|\, \pi_0]
\label{eq:cita_kl}
\end{equation}

where the preference probability is:

\begin{equation}
P^+ = \frac{\exp(\beta \cdot s^+)}{\exp(\beta \cdot s^+) + \exp(\beta \cdot s^-)}
\label{eq:preference_prob}
\end{equation}

with instruction-conditioned scores:
\begin{align}
s^+ &= \log \pi_\theta(Y^+ \mid I, X) \label{eq:score_pos} \\
s^- &= \log \pi_\theta(Y^- \mid I, X) \label{eq:score_neg}
\end{align}

\vspace{-2mm}
\subsection{Hyperparameters}
\label{sec:cita_hyperparams}

\begin{table}[H]
\centering
\small
\begin{tabular}{@{}ll@{}}
\toprule
\textbf{Symbol} & \textbf{Description} \\
\midrule
$\pi_\theta$ & Fine-tuned CITA model \\
\hline
$\pi_0$ & Reference model (pre-CITA checkpoint) \\
\hline
$\beta$ & Temperature (controls preference sharpness) \\
\hline
$\lambda$ & KL weight (\textbf{mandatory}, $> 0$) \\
\bottomrule
\end{tabular}
\caption{CITA-KL hyperparameters. The KL weight $\lambda$ is non-zero by design.}
\label{tab:cita_params}
\vspace{-4mm}
\end{table}

\vspace{-2mm}
\subsection{Training Pipeline}
\label{sec:training_pipeline}

CITA follows a staged training approach, where each stage builds on the previous checkpoint:

\begin{figure}[H]
\centering
\small
\setlength{\tabcolsep}{2pt}
\begin{tabular}{ccccccc}
\textit{\scriptsize Llama-3.1-8B} & & \textit{\scriptsize PKU (chosen)} & & \textit{\scriptsize PKU (pairs)} & & \textit{\scriptsize + Instructions} \\[2pt]
$\downarrow$ & & $\downarrow$ & & $\downarrow$ & & $\downarrow$ \\[2pt]
\fbox{\textbf{Base}} & $\rightarrow$ & \fbox{\textbf{SFT}} & $\rightarrow$ & \fbox{\textbf{DPO}} & $\rightarrow$ & \colorbox{green!20}{\fbox{\textbf{CITA}}} \\[2pt]
& & $\downarrow$ & & $\downarrow$ & & $\downarrow$ \\[2pt]
& & $\mathcal{L}_{\text{SFT}}$ & & $\mathcal{L}_{\text{DPO}}$ & & $\mathcal{L}_{\text{CITA}} + \lambda \cdot \text{KL}$ \\
\end{tabular}
\caption{CITA training pipeline. Each stage uses the previous checkpoint. CITA adds mandatory KL regularization for instruction-conditioned alignment.}
\label{fig:pipeline}
\end{figure}

\vspace{-2mm}
\subsection{Benefits of CITA}
\label{sec:benefits}

\begin{enumerate}[leftmargin=*,itemsep=2pt]
    \item \textbf{Generalization}: Responds appropriately to novel instruction-prompt combinations not seen during training, similar to zero-shot generalization in instruction tuning~\cite{wei2022finetuned,chung2022scaling}.

    \item \textbf{Robustness}: Resists jailbreak attempts~\cite{zou2023universal,wei2023jailbroken} by conditioning behavior on explicit alignment instructions.

    \item \textbf{Calibration}: Avoids over-refusal~\cite{bianchi2023safetytuned} by adjusting safety level based on instruction context.

    \item \textbf{Behavioral Switching}: Dynamically adapts response style (e.g., concise vs. detailed, formal vs. casual) based on instructions, enabling multi-tenant deployment~\cite{wang2024survey}.
\end{enumerate}

\vspace{-3mm}
\subsection{CITA vs. Prior Methods}
\label{sec:comparison}

\begin{table}[H]
\centering
\scriptsize
\begin{tabular}{@{}lccccc@{}}
\toprule
\textbf{Capability} & \textbf{SFT} & \textbf{PPO} & \textbf{GRPO} & \textbf{DPO} & \textbf{CITA} \\
\midrule
No Reward Model & \textcolor{green}{\ding{51}} & \textcolor{red}{\ding{55}} & \textcolor{green}{\ding{51}} & \textcolor{green}{\ding{51}} & \textcolor{green}{\ding{51}} \\
\hline
Preference Pairs & \textcolor{red}{\ding{55}} & \textcolor{red}{\ding{55}} & \textcolor{green}{\ding{51}} & \textcolor{green}{\ding{51}} & \textcolor{green}{\ding{51}} \\
\hline
Online Generation & \textcolor{red}{\ding{55}} & \textcolor{green}{\ding{51}} & \textcolor{green}{\ding{51}} & \textcolor{red}{\ding{55}} & \textcolor{red}{\ding{55}} \\
\hline
Instruction-Conditioned & \textcolor{red}{\ding{55}} & \textcolor{red}{\ding{55}} & \textcolor{red}{\ding{55}} & \textcolor{red}{\ding{55}} & \textcolor{green}{\ding{51}} \\
\hline
Dynamic Policy Switch & \textcolor{red}{\ding{55}} & \textcolor{red}{\ding{55}} & \textcolor{red}{\ding{55}} & \textcolor{red}{\ding{55}} & \textcolor{green}{\ding{51}} \\
\hline
Explicit KL Control & \textcolor{red}{\ding{55}} & \textcolor{green}{\ding{51}} & \textcolor{red}{\ding{55}} & Implicit & \textbf{Mandatory} \\
\bottomrule
\end{tabular}
\caption{Feature comparison across alignment methods. CITA uniquely supports instruction-conditioned behavioral switching with mandatory KL regularization.}
\label{tab:feature_comparison}
\vspace{-4mm}
\end{table}
